%% kapa: exam=1 round=33 year=2022 period=2 form=A subject=appraisal_law subject_ko=감정평가관계법규 seq=15 number=15
%% source: Q-Net 감정평가사 기출문제 (cst00309) artlSeq=5209997 file=raw/1008/5209997/2022년도 제33회 감정평가사 1차시험 2교시(감정평가 관계법규, 회계학).pdf
%% conversion: kordoc 4.24.51 -> tools/gen1.py (auto); review=proofread-sonnet; flags=none
\begin{question}{15}
감정평가 및 감정평가사에 관한 법령상 감정평가에 관한 설명으로 옳지 않은 것은?
\begin{choices}
  \item 금융기관이 대출과 관련하여 토지등의 감정평가를 하려는 경우에는 감정평가법인등에 의뢰하여야 한다.
  \item 감정평가법인등이 해산하거나 폐업하는 경우 시·도지사는 감정평가서의 원본을 발급일부터 5년 동안 보관해야 한다.
  \item 국토교통부장관은 감정평가서가 발급된 후 해당 감정평가가 법률에서 정하는 절차와 방법 등에 따라 타당하게 이루어졌는지를 직권으로 조사할 수 있다.
  \item 최근 3년 이내에 실시한 감정평가 타당성조사 결과 감정평가의 부실이 발생한 분야에 대해서는 우선추출방식의 표본조사가 실시될 수 있다.
  \item 감정평가서에 대한 표본조사는 무작위추출방식으로도 할 수 있다.
\end{choices}
\end{question}
\answer{2}
